\documentclass[11pt,a4paper]{article}
\usepackage[T1]{fontenc}
\usepackage{lmodern,amsmath,amssymb,microtype,booktabs}
\usepackage[margin=25mm]{geometry}
\usepackage[colorlinks=true,urlcolor=blue,linkcolor=blue]{hyperref}
\newcommand{\Om}{\Omega}
\title{The complete dS$_3$ crossing kernel on its heavy-mass peak:\\
phase normalization and an explicit remainder}
\author{Alan Day \and Chengheng\thanks{AI research assistant (OpenAI).}}
\date{Working research note --- 31 October 2026}
\begin{document}\maketitle
\begin{abstract}
After Wilson-duality stabilization, the complete scalar crossing kernel
admits a controlled expansion on the exact collision surface.
We retain its physical CG-period phase, bound the large-parameter
companion series as a whole, and obtain the leading logarithm, constant,
and inverse-mass correction with a compact-transfer uniform remainder.
An explicit finite-mass bound is also given. This is a peak-surface
result for a specified $G_N=0$ model; it neither establishes a
distributional delta limit nor interchanges a physical mass integral
with the heavy-observer limit.
\end{abstract}

\section{Statement, labels and physical scope}
We continue with radius one, $G_N=0$, proper $SO^+(3,1)$, BD scalar
observers and KG matter, a fixed positive observer-label lower bound,
and half-period regulated creation blocks.
The normalized scalar-to-scalar element $U_{00}$ retains the physical
CG-period phase. It is not a separate unitary scalar transform.
The external labels are $(a,q,c,p)$ and internal labels $(b,d)>0$.
Here $M$ is a principal-series label; the static energy is
$\sqrt{1+M^2}$. No energy Jacobian or mass integration is being suppressed.

On the exact collision surface choose
\begin{equation}
 c=M,\quad b=M+u,\quad a=M+u+v,\quad d=M+v,
 \qquad |u|+p\le1,\quad |v|+q\le1 .
 \label{path}
\end{equation}
Define $A_j=1/2+i\alpha_j$ by
\[
 (\alpha_1,\alpha_2,\alpha_3,\alpha_4)
 =\tfrac12(v+q,\ p-u,\ -u-p,\ v-q).
\]
For fixed positive $p,q$ and fixed $u,v$ in this window, we obtain
\begin{equation}
 \boxed{\displaystyle
 U_{00}=
 \frac{\log M}{\pi}
 +\frac{2\psi(1)-\Re\sum_j\psi(A_j)}{2\pi}
 +\frac{u+v}{2\pi M}
 +O\!\left(\frac{\log M}{M^2}\right).}
 \label{main}
\end{equation}
The remainder is uniform on compact admissible transfer sets with
$p,q>0$. All observer labels exceed the fixed lower bound once $M$
is sufficiently large. The proof below is a model calculation built
from the known LPRS kernel and standard special-function identities
\cite{lprs,groenevelt,dlmf}. No originality or full-duality claim is made.

A logarithmically high ridge by itself determines neither the width
nor the integrated weight of a kernel. Equation \eqref{main} is not
used as a delta replacement in any observable.

\newpage
\section{The physical phase and the two series}
Write $h=(1+ib)/2$, $h'=(1+id)/2$ and $h_i=(1+i\nu_i)/2$.
The preceding Wilson reflection identity gives
\[
 U_{00}=C_+XY+\widehat C_-\,\overline{XY},\qquad
 X=\Om(h,h'),\quad Y=\Om(1-h,h'),\quad |\widehat C_-|=|C_+|.
\]
In the present CG-period convention the stronger exact identity is
\begin{equation}
 \widehat C_-=\overline{C_+},\qquad U_{00}=2\Re(C_+XY).
 \label{reality}
\end{equation}
This is an analytic phase identity, not a numerical real projection.
For reproducibility, the period for three all-conjugated legs is
\[
 N(x,y,z)=(4\pi)^{3/2}4^{-3/2-iS/2}
 \frac{\Gamma(\frac12-\frac{iS}{2})
       \prod_{\nu=x,y,z}\Gamma(\frac12-\frac{i(S-2\nu)}2)}
      {\prod_{\nu=x,y,z}\Gamma(1-i\nu)},\quad S=x+y+z .
\]
The physical phase is that of
$N(a,q,b)N(c,p,-b)/[N(c,q,d)N(a,p,-d)]$.
Substitute this and the explicit four-Gamma reflection ratios into
$\widehat C_-C_+$. All powers of four cancel in phase; the remaining
Gamma arguments occur in conjugate pairs with equal integer weights.
Their product is positive. Combined with the previously proved modulus
identity, this proves \eqref{reality}. The full argument/weight list is
saved with an exact symbolic check for six independent real labels.

On \eqref{path}, LPRS (3.38) separates $X$ into
\[
 X=c_A F+c_B G,\quad
 F={}_4F_3(A_1,A_2,A_3,A_4;1+id,1,1-ib;1),
\]
\[
 G={}_4F_3(A_1+ib,\ldots,A_4+ib;
                  1+i(b+d),1+ib,1+ib;1),
\]
\begin{align*}
 c_A&=\frac{\Gamma(ib)\Gamma(1+ib)\Gamma(A_3)\Gamma(A_4)}
              {\Gamma(A_3+ib)\Gamma(A_4+ib)},\\
 c_B&=\frac{\Gamma(-ib)\Gamma(1+id)\Gamma(A_1+ib)\Gamma(A_2+ib)}
              {\Gamma(A_1)\Gamma(A_2)\Gamma(1+i(b+d))}.
\end{align*}
Set $D_0=C_+c_A$ and $E_0=C_+c_B$. Gamma modulus identities imply
\begin{equation}
 D_0=\frac{e^{i\theta(M)}}{4\pi},\qquad
 \theta(M)=\frac{-p^2v+q^2u+u^2v-uv^2+u-v}{4M^2}+O(M^{-3}).
 \label{phase}
\end{equation}
For the modulus, the $\sinh(\pi b)$ in $|C_+|$ cancels that in
$|\Gamma(ib)\Gamma(1+ib)|$; the four remaining cosh factors cancel.
For the phase, insert the displayed period and apply the shifted
Stirling expansion to each Gamma argument $\pm iM+\eta$.
The small-argument Gamma phases cancel in conjugate pairs.
The $M\log M$, $M$, $\log M$, constant and $M^{-1}$ terms all cancel;
the Bernoulli-polynomial term of order $M^{-2}$ gives \eqref{phase}.
The asymptotic continuous phase is chosen to tend to zero.
The shifts range over a compact set, away from the Stirling branch cut,
so the remainder is uniform.

\newpage
\section{Bounding the entire large-parameter companion}
It is not sufficient to keep just the first term of $G$.
Put $A_*=\max_j|\alpha_j|\le1/2$ and write $G=\sum_{n\ge0}g_n$.
For $k=n+1$, since $d>0$,
\begin{align}
 \left|\frac{g_{n+1}}{g_n}\right|
 &\le \frac{[k^2+(b+A_*)^2]^2}
 {k[k^2+b^2]^{3/2}}
 \nonumber\\
 &=\sqrt{1+\frac{b^2}{k^2}}\,
 \left(1+\frac{2bA_*+A_*^2}{k^2+b^2}\right)^2 .
 \label{ratio}
\end{align}
The upper real parts $1/2$ were enlarged to one in this bound.
Using $\log(1+t)\le t$, the sinh product \cite{product}, and
$\sum_{k\ge1}(k^2+b^2)^{-1}\le\pi/(2b)$ gives, for every $n$,
\begin{equation}
 |g_n|\le B_b,\qquad
 B_b=\sqrt{\frac{\sinh\pi b}{\pi b}}\,
       \exp\!\left(2\pi A_*+\frac{\pi A_*^2}{b}\right).
 \label{B}
\end{equation}

To sum this bound, use the preceding tail-certificate polynomial with
upper imaginary parts $b+\alpha_j$ and lower parts $(b+d,b,b)$:
\[
 \mathcal A(x)=\prod_j[(x+\tfrac12)^2+(b+\alpha_j)^2],
 \quad
 \mathcal B(x)=(x+1)^2[(x+1)^2+(b+d)^2][(x+1)^2+b^2]^2,
\]
\[
 P(x)=\mathcal B(x)(x+1)^3-\mathcal A(x)(x+2)^3
     =x^{10}+\sum_{j=0}^{9}p_jx^j.
\]
Let $K_*=\max_{j<10}|p_j|$ and
$N_*=2+\lceil K_*\rceil$.
Every normalized derivative of $P$ is monic with lower coefficients
bounded by $K_*$. For $x>1+K_*$ its lower powers have total absolute
value smaller than the leading power, by the geometric-series bound.
Thus every derivative of $P$ is positive at $N_*$, and
$P(N_*+y)$ has positive coefficients. Consequently
\[
 |g_{n+1}/g_n|\le((n+1)/(n+2))^{3/2}\quad(n\ge N_*).
\]
Split the sum before $N_*$ and use the integral-test tail:
\begin{equation}
 |G|\le\sum_n|g_n|\le (3N_*+3)B_b=:B_G .
 \label{G}
\end{equation}
$N_*$ is an auxiliary analytic threshold, not a loop length used by
our numerical calculation.

On a compact transfer set $K_*=O(M^8)$.
Gamma asymptotics give
$|E_0|=O(M^{-1/2}e^{-\pi M})$ and
$B_b=O(M^{-1/2}e^{\pi M/2})$.
Hence
\begin{equation}
 |E_0G|\le |E_0|B_G
     =O(M^8 e^{-\pi M/2}).
 \label{suppression}
\end{equation}
This deliberately loose exponential estimate is sufficient for
\eqref{main}. It retains the whole companion series rather than
assuming it is close to one.

\newpage
\section{Finite-mass bound and proof of the expansion}
The small-parameter series $F$ shares the coefficients of the collision
digamma series for $Y$. Let $m=\min(b,d)$ and
\[
 q_0=\psi(1-ib)+\psi(1+id)+2\psi(1)-\sum_j\psi(A_j),
\]
\[
 B_M=2+\operatorname{arsinh}b+\operatorname{arsinh}d
                    +4\log2\sum_j|A_j-1|.
\]
The already established coefficient comparison in the window
$|\alpha_j|\le1/2$ gives
\begin{equation}
 |F-1|\le R_F=\frac{\zeta(2)-1}{1+m^2},\qquad
 |Y-q_0|\le R_Y=\frac{(\zeta(2)-1)B_M}{1+m^2}.
 \label{small}
\end{equation}
For example the term estimate follows from
$(x^2-x+1/2)(x+1)\le x^3$ for $x\ge1$; summing the resulting
$(n+1)^{-2}$ bound gives the first inequality.
The digamma difference-series bound supplies $B_M$ for the second.
These are exact-arithmetic analytic bounds.

Combine \eqref{reality} with $X=c_AF+c_BG$:
\[
 U_{00}=\frac1{2\pi}\Re(e^{i\theta}FY)+2\Re(E_0GY).
\]
For the exact finite-mass quantities, an explicit absolute estimate is
\begin{equation}
 \left|U_{00}-\frac{\Re q_0}{2\pi}\right|
 \le \frac{R_Y}{2\pi}
 +\left[\frac{R_F+|e^{i\theta}-1|}{2\pi}
                +2|E_0|B_G\right](|q_0|+R_Y).
 \label{finite}
\end{equation}
The implementation evaluates $|e^{i\theta}-1|=|4\pi D_0-1|$ without
a phase-unwrapping assumption. Its numerical evaluation, unlike the
analytic inequality, is not interval arithmetic.

Uniformly on compact admissible transfers, the estimates are
\[
q_0=O(\log M),\quad R_F=O(M^{-2}),\quad
R_Y=O(\log M/M^2),\quad e^{i\theta}-1=O(M^{-2}).
\]
Equation \eqref{suppression} then reduces \eqref{finite} to
\[
 U_{00}=\frac{\Re q_0}{2\pi}
                  +O(\log M/M^2).
\]
Finally the standard digamma expansion gives
\[
 \Re[\psi(1-ib)+\psi(1+id)]
 =\log b+\log d+O(M^{-2})
 =2\log M+\frac{u+v}{M}+O(M^{-2}).
\]
This proves \eqref{main}, including the specified phase convention.

The proof has not controlled a neighborhood transverse to
$a+c=b+d$, nor an integral over that transverse direction.
A delta limit would require its width, oscillatory tails, normalization
against test functions and the actual physical spectral weights.
The next calculation must supply those ingredients before using
\eqref{main} inside a Bose four-point function.

\newpage
\section{Numerical record and limitations}
Two transfer choices are tested:
$(u,v,p,q)=(0.3,-0.2,0.6,0.6)$ and $(0.2,0.1,0.7,0.8)$.
For each, six masses $M=8,16,32,64,128,256$ are evaluated with the
reflected complete kernel, 128 terms and 100 decimal digits.
All twelve observed discrepancies from $\Re q_0/(2\pi)$ satisfy
\eqref{finite}, allowing separately for the saved series tail.

For the first transfer choice:
\begin{center}
\begin{tabular}{rrr}
\toprule
$M$ & observed discrepancy & explicit bound \eqref{finite}\\
\midrule
8&$2.81\,10^{-3}$&$1.13\,10^{4}$\\
16&$8.55\,10^{-4}$&$5.15$\\
32&$2.53\,10^{-4}$&$3.16\,10^{-3}$\\
64&$7.29\,10^{-5}$&$8.64\,10^{-4}$\\
128&$2.07\,10^{-5}$&$2.35\,10^{-4}$\\
256&$5.77\,10^{-6}$&$6.36\,10^{-5}$\\
\bottomrule
\end{tabular}
\end{center}
The very loose low-mass bounds are retained. The scaled residual
$M^2|U-U_{\rm expansion}|/\log M$ decreases from approximately
$0.0941$ to $0.0711$ for this sample.
Also $M^2\theta$ approaches the analytic coefficient $0.1625$;
the second sample approaches $0.04525$.
These observations check the formulas but do not replace the proof.

The first independent arithmetic audit failed because it imported a
roughly 70-digit value of $\pi$ while requesting 85--90-digit comparisons.
That audit source and failed output remain unchanged.
A separate replay computes $\pi$ from the Machin formula at 180 digits,
with an alternating-series truncation bound below $10^{-160}$.
With unchanged scientific data and unchanged comparison tolerances,
all 191 replay checks pass.
The checks include exact rational positivity at the auxiliary threshold,
finite-bound reconstruction, complex residuals and Gamma-phase
conjugate pairing. This still does not certify arithmetic rounding.

All code, inputs, versions, full complex outputs, failed checks and
repairs are retained with hashes. The earlier 100-digit parameter-derivative
job remains running at this checkpoint and is not counted as passed.
Final project registration and freezing of the recent stages remain pending.

The peak-surface result advances the complete-kernel analysis, while
the off-surface uniform limit, physical wave-packet integration,
Bose free-observer limit, full Fock domains, finite-$G_N$ dynamics,
and microscopic holographic dictionary remain open.

\begin{thebibliography}{9}
\bibitem{lprs} J. Liu et al., \emph{$d$-dimensional SYK, AdS Loops, and
$6j$ Symbols}, \href{https://arxiv.org/html/1808.00612v3}
{arXiv:1808.00612v3}, Eqs.\ (3.35)--(3.38).
\bibitem{groenevelt} W. Groenevelt, \emph{The Wilson function transform},
\href{https://arxiv.org/html/math/0306424v1}{arXiv:math/0306424v1},
Sec.\ 3 and Proposition 4.4.
\bibitem{dlmf} NIST DLMF,
\href{https://dlmf.nist.gov/5.11}{Sec.\ 5.11}, shifted Stirling and
digamma asymptotics.
\bibitem{product} NIST DLMF,
\href{https://dlmf.nist.gov/4.36}{Sec.\ 4.36.1}, sinh infinite product.
Consulted 31 October 2026.
\end{thebibliography}
\end{document}
